\section{Lebesgue Measure}

There are several ways to define the Lebesgue measure, denoted by $\lambda$.
We propose a construction based on outer measures for the one–dimensional case.
In a first part, we construct the measure on $\mathbb{R}$ and then extend the
result to the multidimensional case using the classical outer measure definition.

There is a more natural way to extend the result to higher dimensions using
product measures. However, the approach presented here provides an opportunity
to examine some proofs that are rarely discussed when dealing directly with
multidimensional constructions.

\subsection{The one–dimensional case}

\defn Let us define
\[
	\lambda^*(A) =
	\inf\left\{
	\sum_{n=0}^{+\infty} (b_n-a_n)
	\ :\
	A \subset \bigcup_{n=0}^{+\infty} ]a_n,b_n[
	\right\}
\]
for any $A \subset \mathbb{R}$.

\prop The following properties hold:
\begin{enumerate}
	\item $\lambda^*$ is an outer measure.

	\item $\mathcal{B}(\mathbb{R}) \subset \mathcal{M}(\lambda^*)$.

	\item $\lambda^*(]a,b[)=b-a$.
\end{enumerate}

\rem Although the definition of $\lambda^*$ is given in terms
of countable coverings, we may of course also use finite coverings when
convenient.

Indeed, if a set $A$ is covered by finitely many intervals
\[
	A \subset \bigcup_{n=0}^N ]a_n,b_n[,
\]
then for any $\varepsilon>0$ we can complete this family into a countable
covering by adding intervals of arbitrarily small total length, for instance
$]0,\varepsilon 2^{-n}[$ for $n\ge N+1$. In this way we obtain a countable
covering of $A$ whose total length is
\[
	\sum_{n=0}^N (b_n-a_n) + \sum_{n=N+1}^{+\infty} \varepsilon 2^{-n}
	\le
	\sum_{n=0}^N (b_n-a_n)+\varepsilon.
\]

Since $\varepsilon>0$ is arbitrary, any estimate obtained from a finite
covering may be used in the definition of $\lambda^*$.

\begin{proof}
	Let $A$ be a subset of $\mathbb{R}$.

	\emph{(1)} Let us show that $\lambda^*$ is an outer measure.
	First, let us notice that $\lambda^*(\varnothing)=0$. Indeed,
	\[
	\varnothing \subset ]0,\varepsilon [.
		\]
		So,
		\[
			\lambda^*(\varnothing) \le \varepsilon.
		\]
		It is then clear that $\lambda^*(\varnothing) = 0$.
		Next, by basic properties of the infimum, if $A \subset B$ then
		\[
			\lambda^*(A)\le \lambda^*(B).
		\]
		It remains to prove countable subadditivity.
		Let $(A_n)_{n\in\mathbb{N}}$ be a sequence of subsets of $\mathbb{R}$.
		If
		\[
			\sum_{n=0}^{+\infty}\lambda^*(A_n)=+\infty,
		\]
		then clearly
		\[
			\lambda^*\!\left(\bigcup_{n=0}^{+\infty}A_n\right)
			\le
			\sum_{n=0}^{+\infty}\lambda^*(A_n).
		\]
		Assume now that the series is finite, so that
		$\lambda^*(A_n)<+\infty$ for every $n$.
		Fix $\varepsilon>0$.
		For each $n\in\mathbb{N}$ there exists a sequence of intervals
		$]a_{i,n},b_{i,n}[$ covering $A_n$ such that
		\[
			\sum_{i=0}^{+\infty}(b_{i,n}-a_{i,n})
			\le
			\lambda^*(A_n)+\varepsilon 2^{-n}.
		\]
		The family of intervals $(]a_{i,n},b_{i,n}[)_{i,n\in\mathbb{N}}$
		covers $\displaystyle\bigcup_{n=0}^{+\infty}A_n$.
		Therefore
		\[
			\lambda^*\!\left(\bigcup_{n=0}^{+\infty}A_n\right)
			\le
			\sum_{n=0}^{+\infty}\sum_{i=0}^{+\infty}(b_{i,n}-a_{i,n})
			\le
			\sum_{n=0}^{+\infty}\lambda^*(A_n)+\varepsilon.
		\]
		Since $\varepsilon>0$ is arbitrary, we obtain
		\[
			\lambda^*\!\left(\bigcup_{n=0}^{+\infty}A_n\right)
			\le
			\sum_{n=0}^{+\infty}\lambda^*(A_n).
		\]
		This proves countable subadditivity and therefore $\lambda^*$ is an outer measure.

		\medskip

		\emph{(2)} Let us show that $\mathcal{B}(\mathbb{R}) \subset \mathcal{M}(\lambda^*)$. Recall that $\mathcal{B}(\mathbb{R})=\sigma(]-\infty,a]\,:\,a\in\mathbb{R})$.
		It is therefore sufficient to show that every interval of the form
	$]-\infty,a]$ is $\aM(\lambda^*)$–measurable.
		Let $a\in\mathbb{R}$ and let $(]a_n,b_n[)_{n\in\mathbb{N}}$ be a countable
		family of intervals covering a set $A\subset\mathbb{R}$.
		Then we may decompose the sum of their lengths as
		\begin{align}
			\sum_{n=0}^{+\infty}(b_n-a_n)
			 & =
			\sum_{\substack{n\ge0 \\ b_n\le a}}(b_n-a_n)
			+
			\sum_{\substack{n\ge0 \\ a_n<a<b_n}}(b_n-a_n)
			+
			\sum_{\substack{n\ge0 \\ a\le a_n}}(b_n-a_n).
		\end{align}
		Let $\varepsilon>0$ and define
		\[
			c_n=a-\varepsilon2^{-n}, \qquad d_n=a+\varepsilon2^{-n}.
		\]
		Then
		\[
			4\varepsilon
			+
			\sum_{\substack{n\ge0 \\ a_n<a<b_n}}(b_n-a_n)
			\ge
			\sum_{\substack{n\ge0 \\ a_n<a<b_n}}(b_n-c_n)
			+
			\sum_{\substack{n\ge0 \\ a_n<a<b_n}}(d_n-a_n).
		\]
		Combining these inequalities yields
		\[
			4\varepsilon
			+
			\sum_{n=0}^{+\infty}(b_n-a_n)
			\ge
			\lambda^*(A\cap]-\infty,a])
			+
			\lambda^*(A\cap]a,+\infty[).
		\]
		Taking the infimum over all coverings of $A$ gives
		\[
			4\varepsilon+\lambda^*(A)
			\ge
			\lambda^*(A\cap]-\infty,a])
			+
			\lambda^*(A\cap]a,+\infty[).
		\]
		Since $\varepsilon>0$ is arbitrary, we obtain
		\[
			\lambda^*(A)
			\ge
			\lambda^*(A\cap]-\infty,a])
			+
			\lambda^*(A\cap]a,+\infty[).
		\]
		Hence $]-\infty,a]$ is $\aM(\lambda^*)$–measurable. Since the Borel
	$\sigma$–algebra is generated by such intervals, we conclude that
		\[
			\mathcal{B}(\mathbb{R})\subset\mathcal{M}(\lambda^*).
		\]

		\medskip

		\emph{(3)} Let us prove that $\lambda^*(]a,b[]) = b-a$.
		The interval $]a,b[$ itself covers $]a,b[$, then
		\[
			\lambda^*(]a,b[)\le b-a.
		\]
		To prove the other inequality, choose $a<c<d<b$. We have
		\[
			\lambda^*([c,d])\le \lambda^*(]a,b[).
		\]
		Let $(]a_n,b_n[))_{n\in\mathbb{N}}$ be a covering of $]a,b[$. Then it is also a covering of $[c,d]$, which is compact. Hence there exists $N\in\mathbb{N}$ such that the finite family $(]a_n,b_n[))_{0\le n\le N}$ covers $[c,d]$.
		Now let us denote
		\[
			D_m=\left\{L_k^m=[k2^{-m},(k+1)2^{-m}[\,:\,k\in\mathbb{Z}\right\}.
		\]
		Note that $|L_k^m|$ denotes the length of the interval $L_k^m$, hence
	$|L_k^m|=2^{-m}$.
		For any $n\in\{0,\dots,N\}$, let $A_{n,m}$ be the set of indices $k$
		such that
		\[
			L_k^m\cap ]a_n,b_n[\neq\varnothing .
		\]
		Observe that the intervals $(L_k^m)_{k\in A_{n,m}}$ are contiguous.
		Indeed, since $]a_n,b_n[$ is an interval, if $L_{k_1}^m$ and
	$L_{k_2}^m$ intersect $]a_n,b_n[$ with $k_1<k_2$, then every
	$L_k^m$ with $k_1\le k\le k_2$ also intersects $]a_n,b_n[$.
		Thus the union of these intervals is itself an interval.
		Moreover,
		\[
			]a_n,b_n[
			\subset
			\bigcup_{k\in A_{n,m}} L_k^m .
		\]
		This union may extend slightly outside $]a_n,b_n[$ near the two
		endpoints $a_n$ and $b_n$. At most two dyadic intervals may cross these
		boundary points, and each has length $2^{-m}$. Therefore
		\[
			\sum_{k\in A_{n,m}} |L_k^m|
			\le
			(b_n-a_n)+2\cdot 2^{-m}.
		\]
		Rearranging gives
		\[
			b_n-a_n
			\ge
			\sum_{k\in A_{n,m}} |L_k^m| - 2\cdot 2^{-m}.
		\]
		Summing over $n=0,\dots,N$ yields
		\[
			\sum_{n=0}^{N}(b_n-a_n)
			\ge
			\sum_{n=0}^{N}\sum_{k\in A_{n,m}} |L_k^m| - N2^{-m+1}.
		\]
		It is clear that the chosen intervals $L_k^m$ cover the union of the
		intervals $]a_n,b_n[$, hence they also cover $[c,d]$.
		Let us denote by $B_m$ the set of indices $k$ corresponding to the chosen
		intervals $L_k^m$. Then
		\[
			[c,d] \subset \bigcup_{k\in B_m} L_k^m
			= \bigcup_{n=0}^{N} \bigcup_{k \in A_{n,m}} L_k^m .
		\]
		From this family we can extract a subset $C_m\subset B_m$ such that the
		intervals $(L_k^m)_{k\in C_m}$ are contiguous and still cover $[c,d]$.
		Since $|L_k^m|$ denotes the length of $L_k^m$, and the intervals
	$(L_k^m)_{k\in C_m}$ form a contiguous family covering $[c,d]$, we have
		\[
			\sum_{k\in C_m} |L_k^m| \ge d-c.
		\]
		Moreover,
		\[
			\sum_{n=0}^{N}\sum_{k\in A_{n,m}} |L_k^m|
			\ge
			\sum_{k\in B_m} |L_k^m|
			\ge
			\sum_{k\in C_m} |L_k^m|.
		\]
		Therefore
		\[
			\sum_{n=0}^{N}\sum_{k\in A_{n,m}} |L_k^m|
			\ge d-c.
		\]
		Combining the previous inequalities, we obtain
		\[
			\sum_{n=0}^{N}(b_n-a_n)\ge d-c-N2^{-m+1}.
		\]
		Letting $m\to+\infty$, we get
		\[
			\sum_{n=0}^{N}(b_n-a_n)\ge d-c.
		\]
		Since
		\[
			\sum_{n=0}^{+\infty}(b_n-a_n)\ge \sum_{n=0}^{N}(b_n-a_n),
		\]
		it follows that every covering $(]a_n,b_n[)_{n\in\mathbb{N}}$ of $]a,b[$ satisfies
		\[
			\sum_{n=0}^{+\infty}(b_n-a_n)\ge d-c.
		\]
		Taking the infimum over all such coverings yields
		\[
			\lambda^*(]a,b[)\ge d-c.
		\]
		Since this holds for every $a<c<d<b$, we conclude by letting $c\to a$ and $d\to b$ that
	\[
		\lambda^*(]a,b[)\ge b-a.
	\]
	Together with the reverse inequality, this proves that
	\[
		\lambda^*(]a,b[)=b-a.
	\]
\end{proof}

\subsection{The multidimensional case}
